% =====================================================================
%  Target label: capitalized net occupant wealth (W). Standalone Data
%  subsection -- \input inside Section \ref{sec:acs}. Self-contained:
%  needs only the abstract, introduction, and the methodology cross-
%  references. Notation matches the rest of the paper: Z (relative
%  score), \tilde{SE} (relative SE), S (structural-break statistic).
%
%  Bib keys: poterba1984, himmelberg2005, jorda2019rate,
%  ahlfeldt2015economics (entries at end of file; the last is already in
%  references.bib).
% =====================================================================

\subsection{Target Label: Capitalized Net Occupant Wealth}
\label{sec:wealth}

The label that supervises the model is the \emph{permanent wealth} capitalized into a
tract's built environment, not a single year's income flow. This follows directly from the
spatial-equilibrium logic of Section~\ref{sec:methodology}: the physical housing stock
reflects the expected present value of the marginal occupant's human capital and accumulated
assets, so the object the imagery encodes is a stock, not a flow. We measure that stock with
a tract-level \textbf{net occupant-wealth index} $W_{i,t}$ that capitalizes the three
components a sorting model identifies---human capital, financial assets, and net housing
equity. Two features of household balance sheets make this index a sharper signal than a
contemporaneous income measure. Tracts with equal income but different ownership rates hold
very different wealth, and two equally valued housing stocks represent very different owner
wealth when one is mortgaged and the other owned free and clear. The index addresses both.

\paragraph{Construction.}
Let $\mathrm{pop}_{i,t}$ denote tract population (ACS B01001). Each component is expressed per
capita, so the index compares tracts of different tenure composition on equal footing.

\emph{(i) Human capital.} Per-capita labor earnings (B20003: wages and self-employment
income, which exclude interest, dividends, and rent, and so do not double-count the asset and
housing terms) are capitalized over a finite working horizon as an ordinary annuity,
\begin{equation}
\mathrm{HC}_{i,t}(r) \;=\; \frac{\mathrm{earn}_{i,t}}{\mathrm{pop}_{i,t}}
\;\times\; a(r, N_{i,t}),
\qquad
a(r,N) \;=\; \frac{1-(1+r)^{-N}}{r},
\label{eq:hc}
\end{equation}
where $r$ is the real discount rate and $N_{i,t}$ is the tract's expected remaining working
years, the population-weighted distance to age 65 across the $16$--$64$ brackets of B01001. A
finite horizon, rather than an infinite annuity, reflects that labor income ends at
retirement: a retiree-heavy tract earns a short horizon and leans on its assets.

\emph{(ii) Financial assets.} Per-capita capital income (B19064: interest, dividends, and net
rental income) is capitalized as a perpetuity at the gross asset yield $r_k$, since financial
wealth does not expire at retirement,
\begin{equation}
\mathrm{CAP}_{i,t}(r_k) \;=\; \frac{1}{r_k}\,\frac{\mathrm{capinc}_{i,t}}{\mathrm{pop}_{i,t}}.
\label{eq:cap}
\end{equation}

\emph{(iii) Net housing equity.} Per-capita aggregate owner-occupied home value (B25082) is
de-levered by the mortgaged share $m_{i,t}$ of owners (B25081) at the prevailing
loan-to-value ratio,
\begin{equation}
\mathrm{EQ}_{i,t} \;=\; \frac{\mathrm{ownval}_{i,t}}{\mathrm{pop}_{i,t}}
\;\big(\underbrace{1 - m_{i,t}\,\mathrm{LTV}_t}_{\text{free-clear}\cdot 1\,+\,\text{mortgaged}\cdot(1-\mathrm{LTV}_t)}\big).
\label{eq:eq}
\end{equation}
The aggregate ratio $\mathrm{LTV}_t = 1 - s_t$ uses the Federal Reserve Financial Accounts
series for owners' equity as a share $s_t$ of household real estate (FRED \texttt{HOEREPHRE}),
averaged over the five-year window of each ACS vintage, so the leverage adjustment moves with
the housing cycle. The index sums the three terms,
\begin{equation}
W_{i,t}(r,r_k) \;=\; \mathrm{HC}_{i,t}(r) \;+\; \mathrm{CAP}_{i,t}(r_k) \;+\; \mathrm{EQ}_{i,t}.
\label{eq:w2}
\end{equation}

\paragraph{Rate parameters.}
We set the real discount rate to $r=0.05$ and the gross asset yield to $r_k=0.045$. A real
capitalization rate near five percent is the standard convention in the user-cost-of-housing
literature, where a durable structure is priced as the discounted stream of its service flow
\citep{poterba1984,himmelberg2005}, and it is the rate at which durable locational amenities
capitalize into the built environment in the quantitative spatial-equilibrium models that
motivate our design \citep{ahlfeldt2015economics}. The gross asset yield $r_k=0.045$ sits
within the long-run real-return range for a diversified portfolio of risky and safe assets
documented over more than a century of cross-country data \citep{jorda2019rate}. Neither
choice meaningfully affects the labels the model trains on. Across the grid
$r\in\{0.02,0.03,0.05,0.07\}$ the relative score defined below is at least $0.995$
Spearman-correlated with its $r=0.05$ counterpart, and across $r_k\in\{0.040,0.045,0.050\}$
the ranking is invariant to three decimal places (Appendix~\ref{app:label-robustness}). We
therefore fix one configuration rather than tuning the rates.

\paragraph{Relative cross-sectional score.}
We strip out each city's level and dispersion by standardizing the logged index within city
and year. Let $\mu^{W}_{c,t}$ and $\sigma^{W}_{c,t}$ be the mean and standard deviation of
$\ln W_{i,t}$ across tracts in city (CBSA) $c$ in year $t$. The label fed to the Spatial
Margin Ranking loss is
\begin{equation}
L_{i,t} \;=\; Z^{W}_{i,t} \;=\; \frac{\ln W_{i,t} - \mu^{W}_{c(i),t}}{\sigma^{W}_{c(i),t}}.
\label{eq:zw}
\end{equation}
This normalization makes the label invariant to city-wide nominal drift---inflation and
aggregate housing appreciation---so the model learns only relative structural position.

\paragraph{Standard errors and structural change, 2014--2024.}
$W_{i,t}$ combines several ACS tables whose sampling errors are correlated, so the delta
method understates its variance and an independent Monte-Carlo propagation would ignore the
cross-table covariance. We use the Census Bureau's Variance Replicate Estimate (VRE) tables,
which provide $R=80$ successive-difference replicates of each component. We recompute the full
index \eqref{eq:w2} and its within-city $z$-score \eqref{eq:zw} on every replicate $\rho$ and
take the successive-difference variance of the relative score,
\begin{equation}
\widetilde{\mathrm{SE}}_{i,t}
\;=\; \sqrt{\;\frac{4}{R}\sum_{\rho=1}^{R}\big(Z^{W,(\rho)}_{i,t} - Z^{W,(0)}_{i,t}\big)^{2}\,},
\label{eq:sdr}
\end{equation}
where $Z^{W,(0)}_{i,t}$ is the point estimate. Re-evaluating the same nonlinear index on each
replicate carries the correlation among tables through to the final standard error. The VRE
tables begin in 2014, so the wealth label's standard errors---and its structural-change
test---span 2014 to 2024.

To separate true structural change from sampling noise we apply a two-tailed $z$-test of the
first against the last covered year,
\begin{equation}
S_{i,(2024,2014)} \;=\; \frac{\big|Z^{W}_{i,2024} - Z^{W}_{i,2014}\big|}
{\sqrt{\widetilde{\mathrm{SE}}_{i,2024}^{2} + \widetilde{\mathrm{SE}}_{i,2014}^{2}}},
\label{eq:sbreak-w}
\end{equation}
and flag a significant move with $I^{\,\mathrm{valid}}_{i}=\mathbb{1}\{p(S_{i,(2024,2014)})<0.05\}$.
The classification is robust to the discount rate: it agrees across all four values of $r$ for
$95.1\%$ of tracts (Appendix~\ref{app:label-robustness}).

\paragraph{Relation to the deployment target.}
The wealth label is tightly rank-aligned with per-capita income---Spearman $\rho = 0.97$ at
$r=0.05$, pooled across the $5.4\times10^{4}$ tracts of all large U.S.\ metros. This alignment
is what lets the deployment step quantile-map predicted ranks onto a target city's income
distribution (Section~\ref{sec:methodology}), where income is the one wealth proxy available
in every city, with negligible reordering---while the supervision itself keeps the
permanent-wealth interpretation the theory requires.

% ---------------------------------------------------------------------
% Bib entries (move to references.bib; ahlfeldt2015economics already present):
%
% @article{poterba1984,
%   author  = {Poterba, James M.},
%   title   = {Tax Subsidies to Owner-Occupied Housing: An Asset-Market Approach},
%   journal = {The Quarterly Journal of Economics}, volume = {99}, number = {4},
%   pages   = {729--752}, year = {1984}}
%
% @article{himmelberg2005,
%   author  = {Himmelberg, Charles and Mayer, Christopher and Sinai, Todd},
%   title   = {Assessing High House Prices: Bubbles, Fundamentals and Misperceptions},
%   journal = {Journal of Economic Perspectives}, volume = {19}, number = {4},
%   pages   = {67--92}, year = {2005}}
%
% @article{jorda2019rate,
%   author  = {Jord\`a, \`Oscar and Knoll, Katharina and Kuvshinov, Dmitry and
%              Schularick, Moritz and Taylor, Alan M.},
%   title   = {The Rate of Return on Everything, 1870--2015},
%   journal = {The Quarterly Journal of Economics}, volume = {134}, number = {3},
%   pages   = {1225--1298}, year = {2019}}
% ---------------------------------------------------------------------
